\clearpage\thispagestyle{empty}
\pdfbookmark[0]{Quellen zu Komet (8. April 1815)}{comet1815apr-sources}
\begin{center}\textbf{Quellen zu Komet (8. April 1815)}\end{center}
\noindent Diese KI-Übertragung ins heutige Deutsch folgt dem Bericht mit der Überschrift Comet in den \href{https://gdz.sub.uni-goettingen.de/id/PPN236007467}{\textit{Werken}, Band~VI} von Gauß (Göttingen, 1874), S.~382--383. Er erschien unter Göttingen in den \textit{Göttingischen gelehrten Anzeigen}, Stück~55 (8. April 1815), S.~\href{https://api.digitale-sammlungen.de/iiif/image/v2/bsb10538686_00545/full/1600,/0/default.jpg}{537}, \href{https://api.digitale-sammlungen.de/iiif/image/v2/bsb10538686_00546/full/1600,/0/default.jpg}{538}. Die Links führen zum digitalisierten Exemplar der Bayerischen Staatsbibliothek.
\par\smallskip\noindent Die Werkausgabe gibt als Länge des Perihels $246^\circ\,7'\,2''$ an, der Erstdruck $146^\circ\,7'\,2''$. In beiden steht Enke, und beide geben für die Rektaszension am 30. März 59 13 3 an. Das Sternbild Auriga heißt in der Werkausgabe Fahrmann, im Erstdruck Fuhrmann. KI-Übertragung und KI-Satz.
\clearpage
